\section{Additional Implementation Details}
\label{app:additional-implementation-details}

\subsection{Parameter Sharing}
\label{app:parameter-sharing}

In our homogeneous-objective experiments, the objectives are symmetric: each local policy solves the same type of reachability task, but for a different assigned target. We therefore share the actor-critic network parameters across local policies. This reduces the number of learned parameters and lets the shared network receive PPO updates from all local-policy objectives.

Parameter sharing does not make the local policies indistinguishable. For each local policy \(i\), the policy input contains both a pooled embedding of the full set of objective vectors and the objective-specific vector \(z_i\) for that policy's assigned target. Thus two policy copies share network weights, but are conditioned on different \(z_i\)'s and receive different local rewards. In Cat Feeder, \(z_i\) encodes the assigned feeding request, including its position and remaining lifetime; in the Atari environments, \(z_i\) encodes the assigned enemy object exposed by the object-centric observation.

\subsection{Urgency-Weighted Scalarization}
\label{app:urgency-scalarization}
Let $L$ denote the target expiry interval (\texttt{target\_expiry\_steps}), distinct from the episode horizon $H$, and let $c_{i,t}$ be target $i$'s age counter in the state before action $t$.
Define its clamped remaining lifetime, urgency, and normalized weight by
\begin{align}
  \ell_{i,t} &= \max\{L-c_{i,t},1\}, &
  q_{i,t} &= \frac{\mathbf{1}[i\text{ active at }t]}{\ell_{i,t}}, \\
  w_{i,t} &=
  \begin{cases}
    q_{i,t}/\sum_j q_{j,t}, & \sum_j q_{j,t}>0,\\
    0, & \text{otherwise}.
  \end{cases}
\end{align}
The weights are computed before movement, counter increments, expiration, or respawn.
Active weights sum to one, reached static targets receive zero weight, and there is no additional normalization by the number of targets.
For example, two active targets with remaining lifetimes 1 and 4 receive weights $0.8$ and $0.2$.
The scalar reward is
\begin{equation}
  r_t = \sum_i w_{i,t} r^{\mathrm{obj}}_{i,t}
        -\lambda_B\mathbf{1}[\text{battery depletion at }t],
\end{equation}
where $r^{\mathrm{obj}}_{i,t}$ includes distance-improvement shaping, priority-weighted completion reward, and the expiry penalty.
The depletion term applies only in battery-enabled environments and is not multiplied by target urgency weights.

\subsection{Heterogeneous Cat Feeder Implementation}
\label{app:heterogeneous-implementation}
\paragraph{Battery and stations.}
The reported moving-station comparison uses battery capacity 50.
The three stations start at $(0,0)$, $(15,15)$, and $(29,29)$ and follow independent bounded persistent random walks, moving every 15 steps with direction-change probability $0.1$.
A successful one-cell movement consumes one battery unit; boundary or no-movement steps consume none.
Occupying a charging station immediately restores the battery to capacity.
Reaching zero away from a station causes an immediate tow to the nearest station and a full refill; we count such events as battery depletions.
A capacity-40 fixed-station reference configuration is a different benchmark and is not the configuration used for the reported moving-station results.

\paragraph{Effective bids and payments.}
Feeders and charger both choose raw categorical bids in $\{0,\ldots,6\}$.
Let $B_t$ be the current battery level and $e_t$ the energy required to reach the nearest station.
The charger is forced to bid zero unless $B_t\leq e_t+8$.
For a permitted raw charger bid $b_{\mathrm{chg},t}$, the bid used to determine the winner is
\begin{equation}
  b^{\mathrm{eff}}_{\mathrm{chg},t}
  = b_{\mathrm{chg},t}
    +2\mathbf{1}[b_{\mathrm{chg},t}>0\ \land\ B_t\leq25].
\end{equation}
Thus the effective charger ceiling is 8 when the boost applies, versus 6 for feeders.
The All-Pay payment uses the raw bid, not the boosted effective bid.
This battery-dependent priority rule is an implementation extension to the common-cap auction analyzed in the theory.

\paragraph{Parameter sharing.}
The eight feeder controllers share one actor--critic network, with observations reordered according to the assigned target identity.
The charger uses a completely separate actor--critic with separate direction and bid actor trunks in the moving-station experiment.
No parameters are shared between feeders and charger.

\paragraph{Centralized baseline.}
The unassisted centralized PPO baseline in Table~\ref{tab:heterogeneous-cat-feeder} directly chooses body movement without bidding, an action window, or an execution-time action override.
Let $d_{i,t}$ be distance to target $i$, $p_i\in\{1,2,3,4\}$ its feeding priority, and $D_t$ the battery-depletion indicator.
Its unweighted scalar reward is
\begin{equation}
\begin{split}
 r_t={}&\sum_i\bigl[0.6(d_{i,t}-d_{i,t+1})
       +20p_i\mathbf{1}[i\text{ reached at }t]\\
       &\hspace{3em}-50\mathbf{1}[i\text{ expired at }t]\bigr]-10D_t.
\end{split}
\end{equation}
An assisted centralized variant adds reserve-gated progress shaping toward a station, a refill bonus proportional to the restored battery fraction, and an additional depletion penalty of $100$.
It also has no execution-time action override; these auxiliary terms are not part of the unassisted baseline reported here.

\subsection{Evaluation Aggregation}
\label{app:evaluation-aggregation}
For Table~\ref{tab:all_methods_comp}, PPO results first average each seed's last five available evaluation checkpoints, then report the mean and sample standard deviation across seed averages.
The local-observation ablation (\texttt{visible\_targets=0}, called blind in the discussion) uses the same aggregation as the other PPO variants.
Scaling evaluations use a different checkpoint-selection procedure: a single periodic checkpoint is selected by the largest \texttt{avg\_avg\_performance}, rather than averaging the final five checkpoints.
This selection uses evaluation performance and should be distinguished from evaluating a prespecified final checkpoint on an independent test set.
The scaling plot displays eight times the mean per-episode average performance, which is distinct from the minimum-completion statistic defined in Appendix~\ref{app:scaling-coverage}.
